\documentclass[a4paper,11pt]{article}
\usepackage[utf8]{inputenc}
\usepackage[T1]{fontenc}
\usepackage[french]{babel}
\usepackage[left=2cm,right=2cm,top=2cm,bottom=2cm]{geometry}
\usepackage{xcolor}
\usepackage{hyperref}
\usepackage{fontawesome5}
\usepackage{parskip}

% --- Couleurs INRIA ---
\definecolor{primary}{RGB}{227, 0, 82}

\begin{document}

% --- EXPÉDITEUR (haut gauche) ---
\noindent
\begin{minipage}[t]{0.5\textwidth}
    \textbf{Loïc Aron MBASSI EWOLO}\\
    Élève-Ingénieur ENSEA\\[3pt]
    \faMapMarker\ Cergy, France\\
    \faPhone\ (+33) 7 59 52 33 98\\
    \faEnvelope\ \href{mailto:loicaron.mbassiewolo@ensea.fr}{loicaron.mbassiewolo@ensea.fr}
\end{minipage}
\hfill
% --- LIEU ET DATE (haut droite) ---
\begin{minipage}[t]{0.4\textwidth}
    \raggedleft
    Cergy, le 28 janvier 2026
\end{minipage}

\vspace{1.2cm}

% --- DESTINATAIRE (à droite) ---
\hfill
\begin{minipage}[t]{0.5\textwidth}
    \raggedleft
    \textbf{INRIA Saclay -- Île-de-France}\\
    Équipe TRiBE\\
    Palaiseau, France
\end{minipage}

\vspace{1cm}

\noindent\textbf{Objet :} Candidature Stage Recherche -- Large-scale Privacy Evaluation of Generative Mobility Models\\
\textit{Réf : 2026-09748}

\vspace{0.8cm}

Madame, Monsieur,

Actuellement en double diplôme ingénieur à l'\textbf{ENSEA} et à l'\textbf{École Polytechnique de Yaoundé}, je me permets de vous adresser ma candidature pour le stage de recherche sur l'évaluation à grande échelle de la privacy dans les modèles génératifs de mobilité. Je suis disponible dès \textbf{avril 2026} pour une durée de 6 à 7 mois.

J'ai eu l'opportunité de travailler au \textbf{MILA} (Institut Québécois d'IA) en tant qu'assistant de recherche sur le phénomène de Grokking. Cette expérience m'a formé à la rigueur de la recherche : reproduction d'expériences SOTA, implémentation de pipelines expérimentaux sous \textbf{PyTorch}, documentation rigoureuse et veille bibliographique intensive. J'ai appris à formuler des hypothèses, mener des expérimentations systématiques et analyser les résultats de manière critique.

Votre sujet sur l'évaluation des risques de privacy (\textbf{mémorisation}, \textbf{régurgitation}, similarité entre trajectoires réelles et synthétiques) dans les modèles génératifs de mobilité me passionne. Mon expérience avec les systèmes \textbf{RAG} (projet multi-agents agricole) m'a confronté aux problématiques de cohérence et de traçabilité des informations générées. L'étude des techniques de \textbf{watermarking} et d'\textbf{audit} que vous mentionnez, ainsi que l'analyse des compromis \textbf{privacy--utilité}, représentent des approches que je suis impatient d'explorer sur des datasets publics de mobilité.

Sur le plan technique, je dispose d'une solide expérience en \textbf{Python} (NumPy, Pandas, scikit-learn), \textbf{PyTorch}, et je suis familier avec les workflows expérimentaux à grande échelle. Mes projets open source (Laravel API Generator, +30 étoiles GitHub) témoignent de mon souci de \textbf{reproductibilité} : organisation du code, documentation et versioning rigoureux. Je lis et synthétise couramment des papers en \textbf{anglais scientifique}.

La proximité géographique avec l'INRIA Saclay (je suis à Cergy) et le cadre académique de l'équipe TRiBE me permettraient d'approfondir mon expérience en recherche et de contribuer à des benchmarks reproductibles. Je suis également enthousiaste à l'idée de participer à de potentielles collaborations internationales et de contribuer à la rédaction de rapports techniques ou d'articles scientifiques.

Je serais honoré de contribuer à vos travaux sur l'évaluation des modèles génératifs de mobilité.

Je vous prie d'agréer, Madame, Monsieur, l'expression de mes salutations distinguées.

\vspace{0.8cm}

\textbf{Loïc Aron MBASSI EWOLO}\\
\textit{Élève-Ingénieur ENSEA / Polytechnique Yaoundé}

\end{document}
